\label{sec:intro}

The goal of building robots that perceive, reason, and act in unstructured physical environments has long driven embodied AI research. In recent years, the field has converged on a powerful paradigm: Vision-Language-Action (VLA) models that repurpose pretrained vision-language backbones as generalist robot policies. By formulating action generation as conditional token prediction atop internet-scale visual and linguistic representations, VLA models such as RT-2 \cite{RT-2}, OpenVLA \cite{OpenVLA}, and $\pi_0$ \cite{pi0} have demonstrated striking generalization—following novel language instructions, manipulating previously unseen objects, and transferring across robot embodiments with minimal finetuning. These results establish that the semantic understanding accumulated during large-scale vision-language pretraining can be effectively grounded into motor behavior, marking a qualitative advance over earlier task-specific controllers. The breadth of subsequent work building upon this formulation has consolidated VLA models as the dominant paradigm for generalist embodied policy learning.

However, standard VLA models do not explicitly model world dynamics—they learn direct observation-to-action mappings without predicting how the environment changes under intervention\cite{fei2025liberoplus}. This absence of predictive physical reasoning limits their generalization, where anticipating future states is essential.
Equipping embodied policy models with world modeling capabilities thus emerges as a natural direction \cite{D2PO}.
A growing body of recent work has begun integrating world models into the embodied policy pipeline. These approaches leverage predictive models of environment dynamics to provide agents with physical foresight—whether through video prediction as visual planning \cite{unipi,du2023vlp,ko2024avdc,yang2025roboenvision,gu2026saydreamact}, latent dynamics modeling for policy conditioning \cite{hu2024vpp,pai2025mimicvideo,liang2025videopolicy,yan2026svam,LAPA}, or joint state-action generation within unified architectures \cite{CosmosPolicy,DreamZero,LingBotVA,Motus,UWM,PAD,PhysGen}. This emerging direction has rapidly gained momentum, producing a diverse and expanding landscape of methods.

We formalize this family of approaches under the term \textbf{World Action Models (WAMs)}: embodied foundation models that unify predictive state modeling with action generation, targeting a joint distribution $p(o', a \mid o, l)$ over future states and actions rather than actions alone. Existing WAM methods can be broadly organized into two architectural categories: (1) \textit{Cascaded WAM} explicitly factorizes the objective, formally $p(o', a \mid o, l) = p(a \mid o', o, l)p(o' \mid o, l)$, by first synthesizing representations of anticipated future states, from which actions are subsequently derived; and (2) \textit{Joint WAM} directly modeling the joint distribution ($p(o', a \mid o, l)$)
, where state prediction and action generation are co-optimized within a shared representational space (see \autoref{fig:Treemap} for the temporal evolution of these architectures). The integration of world modeling brings stronger physical understanding, improved generalization across novel environments, and the ability to leverage large-scale human video data that lack action annotations—substantially expanding the data foundations available for embodied policy learning.

This survey provides the \textbf{first} systematic and critical analysis of the World Action Model landscape. Our aim is to offer both a conceptual framework for understanding the design space and a practical guide for researchers entering this rapidly evolving field. The organization of this survey (as illustrated in \autoref{fig:survey-sturcture}) is as follows:



\begin{itemize}
    \item \textbf{Definition (\autoref{sec:def}).} We provide a formal definition of World Action Models and disambiguate them from related concepts, including Video Policies, Action-Conditioned World Models, and standard VLA models, clarifying the terminological boundaries that currently fragment the literature.
    
    \item \textbf{Background (\autoref{sec:background}).} We trace the intertwined development of world modeling and action generation from classical model-based reinforcement learning through modern foundation model approaches, situating WAMs within a broader intellectual lineage.
    
    \item \textbf{Architecture (\autoref{sec:arch}).} We categorize existing WAM methods into Cascaded and Joint paradigms, with further subdivision by generation modality, conditioning mechanism, and action decoding strategy, providing a unified framework for comparing approaches across the design space.
    
    \item \textbf{Training Datasets (\autoref{sec:data}).} We analyze the four major data sources fueling WAM development—robot teleoperation, portable human demonstrations, simulation, and internet-scale egocentric video—examining how each source's characteristics shape the capabilities of trained models.
    
    \item \textbf{Evaluation (\autoref{sec:eval}).} We synthesize the emerging evaluation landscape, organized around visual fidelity, physical commonsense, and action plausibility, and identify gaps where current protocols fall short of assessing WAM-relevant capabilities.

    \item \textbf{Open Challenges (\autoref{sec:open}).} We conclude with a discussion of the critical hurdles and future directions for the field, highlighting the path toward more robust and generalizable World Action Models.
\end{itemize}

